% !TeX program = XeLaTeX
% !TeX root = ../pujavidhanam.tex
\newcommand{\acharya}{श्री-व्यासाचार्य}
\newcommand{\acharyah}{श्री-व्यासाचार्याः}
\newcommand{\acharyan}{श्री-व्यासाचार्यान्}
\newcommand{\acharyebhyonamah}{श्री-व्यासाचार्येभ्यो नमः}

\chapt{श्री-व्यास-पूर्णिमा-पूजा}

\input{purvanga/vighneshwara-puja}

\sect{प्रधान-पूजा — श्री-व्यास-पूजा}

\shuklambaradharam

\pranayama

\renewcommand{\ayane}{दक्षिणायने}
\renewcommand{\rtu}{ग्रीष्म}
\renewcommand{\masa}{कटक-आषाढ}
\renewcommand{\paksha}{शुक्ल}
\renewcommand{\tithau}{पौर्णमास्यां}
\renewcommand{\regularSankalpa}{}
\renewcommand{\additionalSankalpa}{
\begin{itemize}
\item उत्तराषाढा-नक्षत्रे धनूराशौ आविर्भू\-तानां श्रीमत्-शङ्कर-विजयेन्द्र-सरस्वती-श्रीपादानां, शतभिषङ्-नक्षत्रे कुम्भ-राशौ आविर्भूतानां श्रीमत्-सत्य-चन्द्रशेखरेन्द्र-सरस्वती-श्रीपादानाम् अस्माकं जगद्गुरूणां दीर्घ-आयुः-आरोग्य-सिद्ध्यर्थं,

\item तैः सङ्कल्पितानां सर्वेषां लोक-क्षेमार्थ-कार्याणां वेद-शास्त्रादि-सम्प्रदाय-पोषण-कार्याणां विविध-क्षेत्र-यात्रायाश्च अविघ्नतया सम्पूर्त्यर्थं

\item कामकोटि-गुरु-परम्परायां कामकोटि-भक्त-जनानाम् अचञ्चल-भावशुद्ध-दृढतर-भक्ति-सिद्ध्यर्थं, परस्पर-ऐकमत्य-सिद्ध्यर्थं

\item भारतीयानां महाजनानां विघ्न-निवृत्ति-पूर्वक-सत्कार्य-प्रवृत्ति-द्वारा ऐहिक-आमुष्मिक-अभ्युदय-प्राप्त्यर्थम्, असत्कार्येभ्यः निवृत्त्यर्थं

\item भारतीयानां सन्ततेः सनातन-सम्प्रदाये श्रद्धा-भक्त्योः अभिवृद्ध्यर्थं

\item सर्वेषां द्विपदां चतुष्पदाम् अन्येषां च प्राणि-वर्गाणाम् आरोग्य-युक्त-सुख-जीवन-अवाप्त्यर्थम्

\item अस्माकं सह-कुटुम्बानां धर्म-अर्थ-काम-मोक्ष-रूप-चतुर्विध-पुरुषार्थ-सिद्ध्यर्थं विवेक-वैराग्य-सिद्ध्यर्थं
\end{itemize}
}
\renewcommand{\prityartham}{\acharya-प्रीत्यर्थं}
\renewcommand{\kaale}{व्यास-पूर्णिमा-महोत्सवे}
\renewcommand{\pujaam}{\acharya-पूजां}

\sankalpa{}

\vighneshvaraYathasthanam

\input{purvanga/aasana-puja}

\input{purvanga/ghanta-puja}

\input{purvanga/kalasha-puja}

\input{purvanga/aatma-puja}

\input{purvanga/pitha-puja}

\input{purvanga/guru-dhyanam}

\dnsub{ध्यानम्}

\fourlineindentedshloka*
{अभ्र-श्यामः पिङ्ग-जटा-बद्ध-कलापः}
{प्रांशुर्दण्डी कृष्णमृग-त्वक्-परिधानः}
{सर्वान् लोकान् पावयमानः कवि-मुख्यः}
{पाराशर्यः पर्व-सुरूपं विवृणोतु}

\twolineshloka
{व्यासं वसिष्ठ-नप्तारं शक्तेः पौत्रमकल्मषम्}
{पराशरात्मजं वन्दे शुक-तातं तपोनिधिम्}

\twolineshloka
{कृष्ण-द्वैपायनं व्यासं सर्व-भूत-हिते रतम्}
{वेदाब्ज-भास्करं वन्दे शमादि-निलयं मुनिम्}

\twolineshloka
{विश्वरूपं च विश्वेशं विश्व-सत्ता-प्रदं शिवम्}
{वेदयोनिमहं वन्दे व्यासं वेदार्थ-सिद्धिदम्}

अस्मिन् चित्रपटे/पुस्तके/कलशे \acharyan{} ध्यायामि। \acharyan{} आवाहयामि।\\
\acharyebhyonamah{}, आसनं समर्पयामि।\\
\acharyebhyonamah{}, स्वागतं व्याहरामि। पूर्णकुम्भं समर्पयामि।\\
\acharyebhyonamah{}, पाद्यं समर्पयामि।\\
\acharyebhyonamah{}, अर्घ्यं समर्पयामि।\\
\acharyebhyonamah{}, आचमनीयं समर्पयामि।\\
\acharyebhyonamah{}, मधुपर्कं समर्पयामि।\\
\acharyebhyonamah{}, स्नपयामि। स्नानानन्तरम् आचमनीयं समर्पयामि।\\
\acharyebhyonamah{}, वस्त्रं समर्पयामि।\\
\acharyebhyonamah{}, यज्ञोपवीतं समर्पयामि।\\
\acharyebhyonamah{}, दिव्यपरिमलगन्धान् धारयामि।\\ गन्धस्योपरि हरिद्राकुङ्कुमं समर्पयामि।\\
\acharyebhyonamah{}, अक्षतान् समर्पयामि। पुष्पैः पूजयामि।

\begin{center}
    \begingroup
        \setlength{\columnseprule}{1pt}
        \let\chapt\sect
            \input{../namavali-manjari/100/Vyasa_108.tex}
    \endgroup
\end{center}



\acharyebhyonamah{}, नानाविधपरिमल\-पत्र\-पुष्पाणि समर्पयामि।\\
\acharyebhyonamah{}, धूपमाघ्रापयामि।\\
\acharyebhyonamah{}, दीपं दर्शयामि।\\
\acharyebhyonamah{}, अमृतं महानैवेद्यं पानीयं च निवेदयामि। निवेदनानन्तरम् आचमनीयं समर्पयामि।\\
\acharyebhyonamah{}, कर्पूरताम्बूलं समर्पयामि।\\
\acharyebhyonamah{}, मङ्गलनीराजनं दर्शयामि।\\
\acharyebhyonamah{}, प्रदक्षिणनमस्कारान् समर्पयामि।

\acharyebhyonamah{}, प्रार्थनाः समर्पयामि।

\fourlineindentedshloka*
{जयति पराशरसूनुः}
{सत्यवतीहृदयनन्दनो व्यासः}
{यस्यास्यकमलगलितं}
{वाङ्मयममृतं जगत् पिबति}

\sect{व्यास-पूजा-चक्र-देवता-स्मरणम्}

\begin{flushleft}
\begin{multicols}{2}
कृष्णाय शुद्धचैतन्याय नमः\\
वासुदेवाय नमः\\
सङ्कर्षणाय नमः\\
प्रद्युम्नाय नमः\\
अनिरुद्धाय नमः\\
ब्रह्मणे नमः\\
सरस्वत्यै नमः \\
सनकाय नमः\\
सनन्दनाय नमः \\
सनातनाय नमः \\
सनत्कुमाराय नमः \\
सनत्सुजाताय नमः\\
नारदाय नमः\\
वेदव्यासाय नमः\\
शुकाय नमः\\
पैलाय नमः\\
वैशम्पायनाय नमः\\
जैमिनये नमः\\
सुमन्तवे नमः\\
द्रविडाचार्येभ्यो नमः\\
गौडपादाचार्येभ्यो नमः \\
गो\-विन्द\-भगवत्पादा\-चार्येभ्यो नमः \\
शङ्कराचार्येभ्यो नमः\\
पद्मपादाचार्येभ्यो नमः\\
सुरेश्वराचार्येभ्यो नमः\\
हस्तामलकाचार्येभ्यो नमः\\
तोटकाचार्येभ्यो नमः\\
संक्षेपकाचार्येभ्यो नमः\\
विवरणाचार्येभ्यो नमः\\
परात्परगुरुभ्यो नमः\\
परमेष्ठिगुरुभ्यो नमः\\
परमगुरुभ्यो नमः\\
गुरुभ्यो नमः \\
अन्येभ्यो ब्रह्मविद्या\-सम्प्रदाय\-कर्तृभ्य आचार्येभ्यो नमः \\
\end{multicols}
\end{flushleft}

\fourlineindentedshloka
{निगमानपि योऽन्वशाच्चतुर्धा}
{व्यधिताष्टादशधाऽपि यः पुराणम्}
{स च सात्यवतेय ईप्सितं मे}
{सकलाम्नायशिरोगुरुर्विधत्ताम्}

\twolineshloka
{शङ्करं शङ्कराचार्यं केशवं बादरायणम्}
{सूत्रभाष्यकृतौ वन्दे भगवन्तौ पुनः पुनः}

\input{../stotra-sangrahah/stotras/guru/Vyasashtakam.tex}
\input{../stotra-sangrahah/stotras/guru/KanchiKamakotiGuruParamparaStava.tex}

\fourlineindentedshloka*
{जय जय शङ्कर हर हर शङ्कर}
{जय जय शङ्कर हर हर शङ्कर}
{काञ्चीशङ्कर कामकोटिशङ्कर}
{हर हर शङ्कर जय जय शङ्कर}

\fourlineindentedshloka*
{कायेन वाचा मनसेन्द्रियैर्वा}
{बुद्‌ध्याऽऽत्मना वा प्रकृतेः स्वभावात्}
{करोमि यद्यत् सकलं परस्मै}
{नारायणायेति समर्पयामि}

अनेन पूजनेन \acharyah{} प्रीयन्ताम्।

\centerline{ॐ तत्सद्ब्रह्मार्पणमस्तु।}

\closesection
